\scanpage{120}
\mspage{59}
où $\sigma$ est un rev. étale donc un iso. local. Donc on voit que si $D_d\unsure{\text{---}|\longrightarrow}D_0$ et $D_{d'}\unsure{\text{---}|\longrightarrow}D_0$ sont des immersions locales (resp. des mono. locaux), il en est de $=$ de $D_{d,d'}\unsure{\text{---}|\longrightarrow}D_{d'}$. De plus, dans le cas lisse, on voit que la codim. de $D_{d,d'}\unsure{\text{---}|\longrightarrow}D_{d'}$ est égale à la diff. des codim de $D_d\unsure{\text{---}|\longrightarrow}D_0$ et $D_{d'}\unsure{\text{---}|\longrightarrow}D_0$. On supposera, bien sûr (dernier axiome ! (?))
\[
(180)\qquad D_d\unsure{\text{---}|\longrightarrow}D_0\quad\text{est de codim }d
\]

\deleted{\unsure{Et}} Et plus précisément encore ceci. Soit \deleted{\unsure{(181)}} $x\in D_d$, d'où $I(x)$ de cardinal $d$ (tel qu'on ait
\[
(181)\qquad(D'_d)_x\simeq\unsure{\mathcal P}(I_x)
\]
)
\deleted{qui} \deleted{considérons les} on a aussi
\[
(182)\qquad\begin{aligned}
I_x&=\text{germes de sections en }x\\
&\qquad\text{du rev. }D_{d,1}\text{ au dessus de }\unsure{D_d}\\
&\simeq(D_{d,1})_x,
\end{aligned}
\]
Considérons alors les composés $(i\in I)$
\diagram{\begin{tikzcd}[ampersand replacement=\&, column sep=large]
D_d\arrow[r,"s_i"]\&D_{d,1}\arrow[r,"\substack{\mathrm{can}\\\unsure{b|a}}"]\&\unsure{D_1|D_d}\deleted{\unsure{=D_d}}\arrow[r]\&X
\end{tikzcd}}
qui sont tous égaux au morph. can. $D_d\xrightarrow{\mathrm{id}}\deleted{\unsure{D_d}}$, d'autre part les \unsure{images} $x_i$ de $x$ dans $D_1$ pr. les $s_i$.
\margin{119}
